\begin{center}\begin{tikzpicture}[font=\small,>=stealth,line join=round]
\draw[gray!40,rounded corners=3pt] (0,0) rectangle (14.2,8.3);\node at (7.1,7.700000000000001) {隔2变为相反高度，隔4才回到原高度};
\begin{scope}[shift={(1.7153846153846155,3.8842857142857143)},x=0.7692307692307693cm,y=1.3142857142857143cm]\draw[->,gray] (-0.8,0)--(7,0) node[right]{$x$};\draw[->,gray] (0,-1.7)--(0,1.8) node[above]{$y$};\draw[blue,thick,solid,domain=-0.4:6.6,samples=180] plot(\x,{sin(deg(pi*\x/2))});\fill[orange!85!black] (1,1) circle(2.2pt);\fill[orange!85!black] (3,-1) circle(2.2pt);\fill[orange!85!black] (5,1) circle(2.2pt);\foreach \v/\h in {1/1,3/-1,5/1}{\draw[gray,dashed] (\v,0)--(\v,\h);\node[below=5pt] at (\v,0){$\v$};}\end{scope}\node[anchor=west,align=left,text width=5.0cm] at (8.4,4.0) {$f=\sin\frac{\pi x}2,\ a=2$\\[9pt]$f(1)=1$\\[9pt]$f(3)=-1$\\[9pt]$f(5)=1$\\[9pt]正周期可取4};
\node[text width=12.6cm,align=center] at (7.1,.65) {蓝线是满足条件的一条示例曲线；证明只用任意输入的两次变号。};
\end{tikzpicture}\end{center}
